\centerline{\bf \Huge Клетчатые путешествия}

\begin{flushright}
        \scriptsize{нет ничего печальнее, чем зверь запертый в клетке}
\end{flushright}

\problem{1}
Центр города представляет из себя квадрат $5 \times 5$ км, состоящий из 25 кварталов размером $1 \times 1$ км, границы которых - улицы, образующие 36 перекрестков. Какое наименьшее количество полицейских необходимо поставить на перекрестках так, чтобы до каждого из перекрестков какойто из полицейских мог бы добраться, проехав на машине не более 2 км?

\problem{2}
Какое наименьшее число прямоугольников $1 \times 2$ клетки нужно закрасить на доске $8 \times 8$ клеток, чтобы любой квадрат $2 \times 2$ клетки содержал по крайней мере одну закрашенную клетку?

\problem{3}
Какое наибольшее число клеток можно закрасить в квадрате $5 \times 5$ так, чтобы в любом квадрате $2 \times 2$ было закрашено не более двух клеток?


\problem{4}
В белом квадрате $7 \times 7$ поочерёдно закрашиваются в чёрный цвет клетки, у которых до покраски было не более одной чёрной вершины. Какое наибольшее количество клеток можно закрасить таким образом?

\problem{5}
Доска размером $4 \times 4$ клетки покрыта несколькими прямоугольниками размером $1 \times 2$ клетки, стороны которых идут по сторонам клеток. При каком наименьшем числе прямоугольников можно гарантировать, что один из прямоугольников можно убрать так, что оставшиеся будут попрежнему покрывать всю доску?

\problem{6}
Какое наибольшее число клеток можно закрасить на доске $6 \times 6$ так, чтобы они не образовывали букву «Г» из четырёх клеток?